\slajdid{10-s030}
\begin{frame}[label=10-s030]
\gusttekst\setlength{\abovedisplayskip}{3pt}\setlength{\belowdisplayskip}{3pt}\setlength{\abovedisplayshortskip}{2pt}\setlength{\belowdisplayshortskip}{2pt}\begin{columns}[T,onlytextwidth]\column{.36\textwidth}
$t\geq t_2$\par\centering\begin{tikzpicture}[x=.67cm,y=.5cm,font=\sitneoznake,>=Latex]\ctikzset{bipoles/length=.55cm}
\draw(-2,1)node[ocirc]{}node[left]{$V_H$}to[R,l=$R_B$](-.3,1)--(-.3,.3);\draw(-.7,.3)--(.1,.3);\draw(-.55,.13)--(-.05,.13);\draw(-.3,.13)--(-.3,-.6)node[ground]{};\node[left]at(-.7,.1){$V_{BE}$};\node at(-.6,.55){$+$};\draw[->](-1.9,1.4)--(-1.3,1.4)node[midway,above,yshift=3pt]{$I_B$};
\draw(1.5,3)--(1.5,2.8)to[R,l=$R_C$](1.5,1)--(3.5,1)node[ocirc]{};\draw(1.2,3)--(1.8,3);\node[above]at(1.5,3){$V_{CC}$};\draw[->](1.15,2.6)--(1.15,2.1)node[midway,left]{$i_{RC}(t)$};\draw(1.5,1)--(.65,1)--(.65,.55);\node[draw,circle,minimum size=4mm,inner sep=0pt](cs)at(.65,.15){};\draw(.65,.55)--(cs.north);\draw[->](cs.north)--(cs.south);\draw(cs.south)--(.65,-.6)node[ground]{};\node[below]at(.65,-1.3){$\beta_F I_B$};\draw(2.5,1)to[C](2.5,-.9)node[ground]{};\draw[->](2.25,.8)--(1.85,.8);\node at(3.5,1.8){$i_C(t)$};\node[below]at(2.5,-1.3){$C_O$};\draw[->](3.5,-.4)--(3.5,.95)node[midway,right]{$v_o(t)$}node[near end,right]{$+$};
\draw(3.5,-.9)node[ground]{}--(3.5,-.5);
\end{tikzpicture}

\[\tau=C_0 R_C\quad v(t_2^+)=V_{OH}\]
\raggedright\fontsize{9}{10}\selectfont Kondenzator vidi samo otpornost $R_C$. Idealne naponske generatore kratko spajamo, znači ulazni generator $V_H$ i $V_{CC}$ i $V_{BE}$ na masu. Zbog toga je $I_B=0$, pa je i $\beta_F I_B=0$ pa onda taj zavisni strujni izvor možemo da smatramo otvorenom vezom.

\column{.62\textwidth}
\[v_0(t)=V_{CC}-R_C i_{RC}(t)=V_{CC}-R_C(\beta_F I_B-i_C(t))\]
\[i_C(\infty)=0\]
\[v_0(\infty)=V_{CC}-R_C i_{RC}(t)=V_{CC}-R_C\beta_F I_B\]
\[v_0(\infty)=V_{CC}-R_C i_{RC}(t)=V_{CC}-\frac{R_C\beta_F}{R_B}(V_H-V_{BE})\]
\[\frac{R_C\beta_F}{R_B}\gg1\]
\[\frac{R_C\beta_F}{R_B}(V_H-V_{BE})\gg V_{CC}\]
\[\textcolor{DEcrvena}{v_0(\infty)\ll0\quad\text{PO MODELU}}\]
\end{columns}
\[v_o(t)=V_{CC}-\frac{R_C\beta_F}{R_B}(V_H-V_{BE})+\left(V_{OH}-\left(V_{CC}-\frac{R_C\beta_F}{R_B}(V_H-V_{BE})\right)\right)e^{-\frac{t-t_2}{\tau}}\]
\[\textcolor{DEcrvena}{t_3\geq} t\geq t_2\]
\centering\textcolor{DEcrvena}{DOK JE TRANZISTOR U AKTIVNOM REŽIMU}
\end{frame}
